\documentclass[10pt,a4paper]{amsart}
\usepackage[margin=19mm]{geometry}
\usepackage{amsmath,amssymb,mathtools}
\usepackage[scr=rsfs,cal=euler]{mathalfa}
\usepackage{xfrac}
\usepackage[unicode,hidelinks,bookmarks=false]{hyperref}
\newcommand{\ie}{i.e.\ }
\newcommand{\cf}{cf.\ }
\newcommand{\resp}{respectively\ }
\newcommand{\Subsection}{\subsection}
\providecommand{\nobreakdash}{\nobreak\hbox{-}}
\setlength{\emergencystretch}{2em}
\multlinegap0pt
\usepackage{amssymb}
\usepackage{mathtools}
\usepackage[shortlabels]{enumitem}
\usepackage{algorithm}\usepackage{algpseudocode}
\usepackage{braket}
\newtheorem{theorem}{Theorem}
\title{Quantum Estimation of Delay Tail Probabilities in Scheduling and Load Balancing}
\author{R. Srikant}
\date{Source: arXiv:2602.09059v1 (CC BY 4.0)}
\begin{document}
\maketitle

\noindent\emph{Source and licence.} Quantum Estimation of Delay Tail Probabilities in Scheduling and Load Balancing. Version arXiv:2602.09059v1 (CC BY 4.0), distributed by arXiv under the Creative Commons Attribution 4.0 International licence, http://creativecommons.org/licenses/by/4.0/, as recorded in the arXiv licence metadata for this exact version and shown on the abstract page https://arxiv.org/abs/2602.09059v1. Full licence notice: \texttt{licenses/arxiv-2602.09059-CC-BY-4.0.txt}.\par\smallskip

\noindent\textit{Editorial scope.} Counted unit: the article's theorem thm:ER-error-bounded (source0.tex lines 322--344). The article's complete original proof of it is delivered with it (source0.tex lines 346--423, one \texttt{proof} environment of 78 lines). No mathematics is written, repaired or re-derived: the statement and the proof are the article's own lines, in the article's own notation, and no retained line is substituted or re-flowed.  Retained uncounted as the source-local context the proof needs: lines 313--316.  Nothing was omitted from any retained span, so the package carries no omission record.  No retained line contains a citation, so the wrapper carries no bibliography and no bibliography entry.  No retained line contains a figure environment, an image-inclusion command or drawing code, so no image is delivered and the package declares no asset dependency.  Editorial formatting only, in addition: an AMS \texttt{amsart} preamble that restates the article's own declarations and macros used by the retained lines, namely the environments \texttt{theorem} on separate counters, together with the standard packages those lines need.  The retained environments print in this document as Theorem 1 in order of appearance, whereas the article prints the same items with the numbering of its own full document.  The complete native source of the article as distributed in the arXiv e-print for this exact version is bundled unchanged as \texttt{sources/194290/source0.tex}.
\begin{equation}
\label{eq:beta-def}
\beta := \frac{2\Delta^2}{(A_{\max}+S_{\max})^2},
\end{equation}

\begin{theorem}
\label{thm:ER-error-bounded}

Fix a truncation horizon $M$. Let $\hat{\mu}_Q$
be an estimate of $E[Y]=E[R_M]/M$ returned by QAE
such that, with probability at least $1-\delta_Q$,
\begin{equation}
\label{eq:qae-err-bounded}
|\hat{\mu}_Q - E[Y]|
\le \varepsilon_Q.
\end{equation}
Then, with probability at least $1-\delta_Q$, the corresponding estimate
$\widehat{E[R]} := M\,\hat{\mu}_Q$ satisfies
\begin{equation}
\label{eq:ER-bound-bounded}
\Bigl|\,\widehat{E[R]} - E[R]\,\Bigr|
\;\le\;
M\varepsilon_Q
\;+\;
\frac{e^{-\beta M}}{1-e^{-\beta}},
\end{equation}
where $\beta$ is given by~\eqref{eq:beta-def}.
\end{theorem}

\begin{proof}
Define the i.i.d.\ increments
\[
X_i := S_i - A_{i+1},\qquad i\ge 0,
\]
so that $X_i\in[-A_{\max},S_{\max}]$ almost surely and
$E[X_i]=-(E[A_1]-E[S_1])=-\Delta$.

Let $\{V_t\}_{t\ge 0}$ be the associated unreflected random walk
\[
V_0 := 0,
\qquad
V_t := \sum_{i=0}^{t-1} X_i , \quad t\ge 1.
\]
For the reflected Lindley recursion $W_{t+1}=(W_t+X_t)^+$ with $W_0=0$, the
workload admits the standard representation
\begin{equation}
\label{eq:reflection-repr}
W_t = V_t - \min_{0\le k\le t} V_k
= \max_{0\le k\le t} \sum_{i=k}^{t-1} X_i ,
\qquad t\ge 0.
\end{equation}

Recall that $\tau:=\inf\{t\ge 1 : W_t=0\}$ is the first return time to zero.
On the event $\{\tau>t\}$ we have $W_1>0,\dots,W_t>0$.  By
\eqref{eq:reflection-repr}, the condition $W_m=0$ is equivalent to
$V_m=\min_{0\le k\le m}V_k$.  Therefore, $\{\tau>t\}$ implies that the
unreflected walk $V_m$ does not attain a new running minimum over
$m=1,\dots,t$, so that
\[
\min_{0\le k\le t} V_k = V_0 = 0.
\]
Consequently, $V_m>0$ for all $m=1,\dots,t$, and in particular $V_t>0$.
Thus,
\[
P(\tau>t)
\;\le\;
P(V_t>0)
=
P\!\left(\sum_{i=0}^{t-1} X_i \ge 0\right)
=
P\!\left(\sum_{i=0}^{t-1} (X_i-E[X_i]) \ge t\Delta\right).
\]

Applying Hoeffding's inequality to the bounded variables
$X_i-E[X_i]\in[-A_{\max}+\Delta,\,S_{\max}+\Delta]$, whose range length is at
most $A_{\max}+S_{\max}$, yields
\[
P(\tau>t)
\le
\exp\!\left(
-\frac{2t^2\Delta^2}{t(A_{\max}+S_{\max})^2}
\right)
=
\exp(-\beta t),
\]
where $\beta$ is defined in~\eqref{eq:beta-def}.

Since $R-R_M=0$ on $\{\tau\le M\}$ and $0\le R-R_M\le \tau-M$ on $\{\tau>M\}$,
it follows that
\[
0\le E[R]-E[R_M]
\le E[(\tau-M)^+]
= \sum_{t=M}^{\infty} P(\tau>t)
\le \sum_{t=M}^{\infty} e^{-\beta t}
= \frac{e^{-\beta M}}{1-e^{-\beta}}.
\]

If a quantum subroutine returns $\hat{\mu}_Q$ satisfying
\eqref{eq:qae-err-bounded}, then
\[
|M\hat{\mu}_Q - E[R_M]|
= M|\hat{\mu}_Q - E[Y]|
\le M\varepsilon_Q.
\]
Combining the bounds and applying the triangle inequality yields
\eqref{eq:ER-bound-bounded}.
\end{proof}

\end{document}
